\begin{examplebox}
	\textbf{\textcolor{CExample}{例 1（同底加减，先合并）}}

	\textbf{题目：} 计算 $\log_3 81+\log_3 27-\log_3 9$。

	\textbf{解：} 三个对数底数相同，先把加法改成乘法、减法改成除法：
	\[
		\begin{aligned}
		\log_3 81+\log_3 27-\log_3 9
		&=\log_3\frac{81\times27}{9}\\
		&=\log_3 243=5.
		\end{aligned}
	\]
	\textbf{为什么对？} 因为 $243=3^5$。\textbf{何时最快？} 同底对数连加连减，先合并往往比逐项计算省事。

	\medskip
	\textbf{\textcolor{CExample}{例 2（根号里有数，指数提前）}}

	\textbf{题目：} 计算 $\log_2\sqrt8$。

	\textbf{解：} 根号可写成 $\frac12$ 次方：
	\[
		\log_2\sqrt8=\log_2 8^{\frac12}
		=\frac12\log_2 8=\frac12\times3=\boxed{\frac32}.
	\]
	\textbf{为什么对？} 因为 $\sqrt8=2^{3/2}$。\textbf{何时最快？} 真数中出现乘方、平方根，先考虑把指数提出。

	\medskip
	\textbf{\textcolor{CExample}{例 3（底数不同，换底约分）}}

	\textbf{题目：} 计算 $\log_2 3\cdot\log_3 4$。

	\textbf{解：} 都换成以 $2$ 为底：
	\[
		\log_2 3\cdot\log_3 4
		=\log_2 3\cdot\frac{\log_2 4}{\log_2 3}
		=\log_2 4=\boxed{2}.
	\]
	\textbf{为什么对？} 换底后，相同的 $\log_2 3$ 可以约去。\textbf{何时最快？} 不同底数的对数相乘，先观察能否换底约分。

	\medskip
	\textbf{\textcolor{CExample}{例 4（对数方程，化成普通方程）}}

	\textbf{题目：} 已知 $x,y>0$，解方程组
	\[
		\begin{cases}
			\log_2 x+\log_2 y=3,\\
			x-y=2.
		\end{cases}
	\]
	\textbf{解：} 将第一式合并：
	\[
		\log_2(xy)=3\quad\Longrightarrow\quad xy=8.
	\]
	由 $y=x-2$，得到 $x(x-2)=8$，即
	\[
		x^2-2x-8=(x-4)(x+2)=0.
	\]
	因为 $x>0$，舍去 $x=-2$，所以
	\[
		\boxed{x=4,\quad y=2}.
	\]
	\textbf{何时最快？} 方程里有同底对数相加，先合并，就可能把对数方程变成学过的普通方程。
\end{examplebox}
